% Generated draft: requires mathematical and source review.
\documentclass[11pt]{article}
\usepackage[margin=25mm]{geometry}
\usepackage{fontspec}
\setmainfont{Latin Modern Roman}
\usepackage{amsmath,amssymb,mathtools,unicode-math}
\setmathfont{Latin Modern Math}
\usepackage{enumitem,needspace,xurl,hyperref}
\setlength{\parindent}{0pt}
\setlength{\parskip}{5pt}
\setlength{\emergencystretch}{3em}
\newcommand{\N}{\mathbb{N}}
\newcommand{\Z}{\mathbb{Z}}
\newcommand{\Q}{\mathbb{Q}}
\newcommand{\R}{\mathbb{R}}
\newcommand{\C}{\mathbb{C}}
\newcommand{\F}{\mathbb{F}}
\newcommand{\A}{\mathbb{A}}
\newcommand{\PP}{\mathbb{P}}
\newcommand{\E}{\mathbb{E}}
\newcommand{\eps}{\varepsilon}
\newcommand{\dd}{\,\mathrm{d}}
\begin{document}
\section*{UnsolvedMath ID: 800253}
\textbf{Source problem code:} catalogue entry
\subsection{Definitions and mathematical statement}
Let \(n\geq 1\), and let \((V,0)\) be the germ at the origin of an irreducible complex analytic variety in \(\C^n\). Choose a representative on a neighborhood of \(0\). For sufficiently small \(\epsilon>0\), write
\[
B_\epsilon=\{z\in\C^n:|z|<\epsilon\},\qquad E_\epsilon=V\cap\R^n\cap B_\epsilon,
\]
where \(|\cdot|\) denotes the Euclidean norm and \(\Im z=(\Im z_1,\ldots,\Im z_n)\). The source displays the extremal function
\[
U_{E_\epsilon}(z):=\sup\bigl\{u(z):u\leq0\text{ on }E_\epsilon,\quad 0\leq u\leq1\text{ on }B_\epsilon\bigr\},\qquad z\in B_\epsilon.
\]
The admissible function class for \(u\) is not specified in the supplied record. In particular, no continuity, subharmonicity, or plurisubharmonicity condition is stated. Thus the intended extremal function is not fully specified; adding such a condition would be an additional hypothesis.

The source asks to classify those germs \((V,0)\) for which there is a constant \(A>0\) such that
\[
(\alpha)\qquad U_{E_\epsilon}(z)\leq A|\Im z|\qquad\text{for every }z\in B_{\epsilon/2}.
\]
The displayed estimate is stated on the entire ambient ball, without a restriction \(z\in V\). The construction is made for any sufficiently small \(\epsilon\), but the source does not explicitly state whether \(A\) may depend on \(\epsilon\) or must be uniform over all sufficiently small radii. These unresolved qualifications must be settled before this becomes a precise analytic classification problem.
\subsection{Short English statement}
Classify irreducible complex analytic germs at \(0\) whose stated extremal function associated with their real points satisfies a linear bound in the distance to \(\R^n\). The supplied formulation omits the admissible function class and leaves the dependence of the bound on the radius unclear.
\subsection{Sources}
\begin{itemize}
\item Supplied local UnsolvedMath registry record: permanent ID \(800253\), source code \texttt{catalogue entry}, title ``Several Complex Variables Problem List --- Problem 52.'' The registry attributes the record to Buzzard, Fornaess, and Krantz, with proposed year \(1995\), and labels it open; that label is supplied metadata, not an independently verified mathematical status.
\item Buzzard, Fornaess, and Krantz, \emph{The Several Complex Variables Problem List} (1995), Problem 52, marked ``(New)'' in the supplied statement. Supplied source link: \url{https://arxiv.org/abs/math/9509202}. The provenance records access on \(2026\text{-}10\text{-}16\) and extraction by \texttt{arxiv\_latex}; these are supplied provenance claims, not access or verification performed for this draft.
\end{itemize}
\end{document}
